\documentclass[a4paper,12pt]{article}
\usepackage{iftex}
\ifPDFTeX
  \usepackage[T2A]{fontenc}
  \usepackage[utf8]{inputenc}
  \usepackage[russian]{babel}
  \usepackage{DejaVuSans}
  \renewcommand{\familydefault}{\sfdefault}
\else
  \usepackage{fontspec}
  \usepackage[russian]{babel}
  \setmainfont{Arimo}
  \setsansfont{Arimo}
  \renewcommand{\familydefault}{\sfdefault}
\fi
\usepackage{amsmath,amssymb,mathtools}
\usepackage{icomma}
\usepackage{geometry}
\geometry{a4paper,left=21mm,right=21mm,top=21mm,bottom=23mm}
\usepackage{microtype}
\usepackage{xcolor}
\definecolor{accent}{HTML}{13767A}
\definecolor{darkaccent}{HTML}{23535B}
\definecolor{muted}{HTML}{53626A}
\definecolor{rulecolor}{HTML}{C4D4D4}
\usepackage{tikz}
\usetikzlibrary{calc,arrows.meta,positioning,shapes.geometric}
\usepackage{booktabs,tabularx}
\usepackage{enumitem}
\usepackage{hyperref}
\hypersetup{hidelinks,pdftitle={Чек-лист. Оценки выражений и линейные неравенства},pdfauthor={Лёвин Артём Александрович}}
\usepackage{fancyhdr}
\usepackage{setspace}
\usepackage{titlesec}
\onehalfspacing
\setlength{\parindent}{0pt}
\setlength{\parskip}{0.36em}
\setlist[enumerate]{leftmargin=*,itemsep=0.65em,topsep=0.45em}
\titleformat{\section}{\LARGE\bfseries\color{darkaccent}}{}{0pt}{}
\titlespacing*{\section}{0pt}{12pt}{8pt}
\titleformat{\subsection}{\large\bfseries\color{accent}}{}{0pt}{}
\titlespacing*{\subsection}{0pt}{11pt}{4pt}
\pagestyle{fancy}
\fancyhf{}
\fancyhead[L]{\small\textcolor{muted}{Математика \enspace / \enspace 10.10.26}}
\fancyhead[R]{\small\textcolor{muted}{Екатерина Скелли}}
\fancyfoot[C]{\small\textcolor{muted}{\thepage}}
\renewcommand{\headrulewidth}{0.25pt}
\renewcommand{\footrulewidth}{0pt}
\setlength{\headheight}{16pt}
\newcommand{\key}[1]{\textbf{\textcolor{darkaccent}{#1}}}
\newcommand{\note}[1]{\par\smallskip{\color{darkaccent}\bfseries Обратите внимание.}\enspace #1\par\smallskip}
\newcommand{\sol}[1]{\textcolor{accent}{\textbf{#1}}}

\begin{document}
\thispagestyle{empty}
\begin{center}
\vspace*{27mm}
{\large\color{muted} МАТЕМАТИКА \textbullet\ ПОВТОРЕНИЕ ПО ЗАНЯТИЮ}\par
\vspace{21mm}
{\fontsize{36}{42}\selectfont\bfseries\color{darkaccent} Чек-лист}\par
\vspace{13mm}
{\fontsize{23}{30}\selectfont\bfseries Оценка выражений\par и линейные неравенства}\par
\vspace{14mm}
{\large\color{muted} Границы величин \textbullet\ знаки неравенств \textbullet\ целые решения}\par
\vspace{24mm}
{\large 10 октября 2026 года}\par
\vfill
{\large\color{darkaccent} Лёвин Артём Александрович}\par
\vspace{2mm}
{\large\color{darkaccent} эксклюзивно для Екатерины Скелли}\par
\vspace{11mm}
\end{center}
\begin{tikzpicture}[remember picture,overlay]
\draw[accent!55,line width=1pt] ([xshift=20mm,yshift=25mm]current page.south west)
 .. controls ([xshift=65mm,yshift=38mm]current page.south west)
 and ([xshift=-68mm,yshift=15mm]current page.south east)
 .. ([xshift=-20mm,yshift=28mm]current page.south east);
\end{tikzpicture}
\clearpage

\section*{01 \quad Оценка выражений: основные приёмы}
\key{Цель:} установить, между какими числами находится значение выражения, даже если точное значение неизвестно.

\subsection*{1. Сохранение порядка и возведение в степень}
Если $a<b$, то при прибавлении одинакового числа знак сохраняется. При умножении или делении на \key{положительное} число знак тоже сохраняется; на \key{отрицательное} --- меняется на противоположный.

При возведении обеих частей в \key{нечётную положительную степень} порядок сохраняется. Например,
\[
0{,}9<a<1{,}1 \quad\Longrightarrow\quad
0{,}729<a^3<1{,}331.
\]
Для чётной степени при наличии отрицательных чисел требуется отдельно учесть положение нуля: правило «возвести границы и оставить порядок» может дать ошибку.

\subsection*{2. Сложение и вычитание промежутков}
Если $a\le x\le b$ и $c\le y\le d$, то
\[
\boxed{a+c\le x+y\le b+d},\qquad
\boxed{a-d\le x-y\le b-c}.
\]
\key{Логика разности:} наименьшее $x-y$ получается при меньшем $x$ и большем $y$; наибольшее --- наоборот.

\subsection*{3. Умножение промежутков}
Для $xy$ вычислите все четыре произведения концов:
\[
ac,\quad ad,\quad bc,\quad bd.
\]
Наименьшее и наибольшее из них дают границы произведения. Это особенно полезно, когда промежутки содержат отрицательные числа.

\key{Пример.} Пусть $-2\le x\le1$ и $-1\le y\le3$. Тогда
\[
\begin{gathered}
(-2)(-1)=2,\quad(-2)\cdot3=-6,\\
1\cdot(-1)=-1,\quad1\cdot3=3;
\qquad \boxed{-6\le xy\le3}.
\end{gathered}
\]
\clearpage

\section*{02 \quad Деление: границы и ограничения}
\subsection*{4. Положительное число делим на переменную}
При $0<a\le t\le b$ и $C>0$ справедливо
\[
\boxed{\frac{C}{b}\le\frac{C}{t}\le\frac{C}{a}}.
\]
\key{Смысл:} чем больше положительный знаменатель, тем меньше частное. Например, при $3{,}5\le t\le5{,}5$:
\[
\frac{120}{5{,}5}\le\frac{120}{t}\le\frac{120}{3{,}5},
\quad\text{то есть}\quad
\frac{240}{11}\le\frac{120}{t}\le\frac{240}{7}.
\]
Если плейлист длится \key{ровно} 120 минут, а длительность каждой композиции от $3{,}5$ до $5{,}5$ минуты, то для \emph{целого} числа композиций $N$ получаем
\[
3{,}5N\le120\le5{,}5N
\quad\Longrightarrow\quad \boxed{22\le N\le34}.
\]

\subsection*{5. Отношение двух величин}
Чтобы оценить $x/y$, убедитесь, что знаменатель \key{ни при каких допустимых значениях не равен нулю}. Если диапазон $y$ целиком по одну сторону от нуля, проверьте четыре частных, полученных из граничных значений $x$ и $y$.

\key{Пример с отрицательными числами.} Если
\[
-0{,}9<u<-0{,}8,\qquad -0{,}9<v<-0{,}8,
\]
то возможные граничные частные равны $1$, $\frac89$, $\frac98$, $1$. Следовательно,
\[
\boxed{\frac89<\frac uv<\frac98}.
\]
Если $u=b/a$, $v=b/c$ и $b\ne0$, то $\frac uv=\frac ca$.

\note{Если промежуток знаменателя содержит нуль, деление при нуле не определено. Более того, при ненулевом числителе частное может не иметь конечных границ. Так, для $-2\le x\le1$, $-1\le y\le3$ нельзя ограничить $x/y$ двумя конечными числами: можно взять $x=1$ и $y$ сколь угодно близким к нулю.}
\clearpage

\section*{03 \quad Линейные неравенства}
\subsection*{6. Как решить неравенство}
Решение сводится к виду $ax\,<\,b$, $ax\,\le\,b$, $ax\,>\,b$ или $ax\,\ge\,b$.
\begin{enumerate}
\item Раскройте скобки и приведите подобные слагаемые.
\item Перенесите слагаемые с $x$ в одну сторону, числа --- в другую.
\item Разделите обе части на коэффициент при $x$. Если он \key{отрицательный}, разверните знак.
\item Запишите ответ промежутком; при необходимости отметьте его на числовой прямой.
\end{enumerate}

\key{Пример 1.}
\[
-2x>14\quad\Longrightarrow\quad\boxed{x<-7},
\qquad x\in(-\infty;-7).
\]

\key{Пример 2.}
\[
\begin{aligned}
2x+3&\ge5x-6,\\
-3x&\ge-9,\\
x&\le3,\qquad x\in(-\infty;3].
\end{aligned}
\]

\subsection*{7. Строгость и запись промежутка}
\begin{tabularx}{\linewidth}{@{}lXX@{}}
\toprule
\textbf{Знак} & \textbf{Числовая прямая} & \textbf{Граница в ответе}\\
\midrule
$<$, $>$ & точка не включена (пустой кружок) & круглая скобка\\
$\le$, $\ge$ & точка включена (закрашена) & квадратная скобка\\
\bottomrule
\end{tabularx}

\note{У $-\infty$ и $+\infty$ всегда круглые скобки. Строгий знак запрещает равенство граничному числу.}

\subsection*{8. Перевод слов в знаки}
\begin{tabularx}{\linewidth}{@{}XX@{}}
\toprule
\textbf{Формулировка} & \textbf{Математическая запись}\\
\midrule
отрицательное / положительное & $A<0$ / $A>0$\\
неотрицательное / неположительное & $A\ge0$ / $A\le0$\\
не больше $B$ / не меньше $B$ & $A\le B$ / $A\ge B$\\
\bottomrule
\end{tabularx}
\clearpage

\section*{Алгоритм \quad Линейное неравенство}
\vspace{4mm}
В схеме $\triangleleft$ означает один из знаков $<$, $\le$, $>$, $\ge$.
\vspace{5mm}
\begin{center}
\begin{tikzpicture}[
 >=Stealth,
 flowstep/.style={draw=rulecolor,rounded corners=2mm,align=center,fill=white,text width=8.3cm,minimum height=1.05cm,font=\small},
 decision/.style={draw=accent,diamond,aspect=2,align=center,inner sep=1.2mm,font=\small,minimum width=2.8cm,minimum height=1.45cm},
 branch/.style={draw=rulecolor,rounded corners=1.8mm,align=center,text width=3.75cm,minimum height=1.05cm,font=\small},
 result/.style={draw=accent,rounded corners=4mm,align=center,text width=2.4cm,minimum height=0.88cm,font=\small\bfseries,color=darkaccent},
 arrow/.style={-{Stealth[length=2.7mm]},line width=.68pt,draw=darkaccent},
 node distance=8mm
]
\node[flowstep] (s) {Раскройте скобки, приведите подобные,\\ избавьтесь от числовых знаменателей};
\node[flowstep,below=of s] (eq) {Получите вид $ax\ \triangleleft\ b$};
\node[decision,below=12mm of eq] (z) {$a=0$?};
\node[decision] (truth) at (-4.25,-6.4) {$0\ \triangleleft\ b$\\ верно?};
\node[decision] (sign) at (4.25,-6.4) {$a>0$?};
\node[result] (r1) at (-6.15,-9.65) {Все $x\in\mathbb R$};
\node[result] (r2) at (-2.2,-9.65) {Решений нет};
\node[branch] (r3) at (2.1,-9.65) {Делите на $a$,\\ знак сохраняйте};
\node[branch] (r4) at (6.15,-9.65) {Делите на $a$,\\ знак меняйте};
\draw[arrow] (s) -- (eq);
\draw[arrow] (eq) -- (z);
\draw[arrow] (z.west) -| node[pos=.27,above,font=\small]{да} (truth.north);
\draw[arrow] (z.east) -| node[pos=.27,above,font=\small]{нет} (sign.north);
\draw[arrow] (truth.south west) -- node[midway,left,font=\small]{да} (r1.north);
\draw[arrow] (truth.south east) -- node[midway,right,font=\small]{нет} (r2.north);
\draw[arrow] (sign.south west) -- node[midway,left,font=\small]{да} (r3.north);
\draw[arrow] (sign.south east) -- node[midway,right,font=\small]{нет} (r4.north);
\end{tikzpicture}
\end{center}
\vspace{5mm}
После деления укажите множество решений. Для вопроса о \emph{целых} решениях оставьте только целые числа из найденного промежутка.
\clearpage

\section*{04 \quad Особые случаи и приложения}
\begin{spacing}{1.12}
\small
\subsection*{9. Когда переменная исчезла}
\key{Верное числовое неравенство} означает, что подходят все действительные $x$:
\[
x+6>x\quad\Longleftrightarrow\quad6>0
\quad\Longrightarrow\quad x\in\mathbb R.
\]
\key{Ложное числовое неравенство} означает отсутствие решений:
\[
x+6<x\quad\Longleftrightarrow\quad6<0
\quad\Longrightarrow\quad S=\varnothing.
\]

\subsection*{10. Дроби в линейном неравенстве}
Если знаменатели --- известные положительные числа, умножьте \key{каждое слагаемое обеих частей} на их наименьшее общее кратное. Например,
\[
\frac{13x-1}{15}-\frac{2x-1}{5}<x-\frac{x-2}{3}.
\]
Умножаем на $15>0$, поэтому знак не меняется:
\[
\begin{aligned}
13x-1-3(2x-1)&<15x-5(x-2),\\
7x+2&<10x+10,\\
-3x&<8,\\
x&>-\frac83.
\end{aligned}
\]
Среди \key{отрицательных целых} решений здесь только $-2$ и $-1$.

\subsection*{11. Квадраты, которые сокращаются}
Иногда неравенство кажется квадратным, но после раскрытия скобок $x^2$ исчезает:
\[
\begin{aligned}
(x+4)^2&\le x^2+5x+4,\\
x^2+8x+16&\le x^2+5x+4,\\
3x&\le-12,\qquad\boxed{x\le-4}.
\end{aligned}
\]

\subsection*{12. Ограничения в текстовой задаче}
Корзинка стоит 150 руб., один тюльпан --- 45 руб., бюджет --- 2000 руб. При $n\in\mathbb Z$, $n\ge0$:
\[
45n+150\le2000\quad\Longrightarrow\quad
n\le\frac{1850}{45}=41\frac19.
\]
Значит, можно купить \key{не более 41 тюльпана}. Дробное число предметов в таком контексте не допускается.
\end{spacing}
\clearpage

\section*{05 \quad Тренировка \textemdash{} Путь А}
\textcolor{muted}{10 заданий от базовой техники к применению. Оформите решение и запишите ответ.}
\begin{enumerate}[label=\textbf{\arabic*.},itemsep=1.0em]
\item Известно, что $0{,}7<a<1{,}1$. Оцените $a^3$.
\item Известно, что $4<t<6$. Укажите промежуток для $\dfrac{24}{t}$.
\item Пусть $-2\le k\le1$, $-1\le n\le3$. Оцените $k+n$ и $k-n$.
\item При тех же условиях найдите границы $kn$. Можно ли указать \emph{конечные} нижнюю и верхнюю границы $k/n$ для всех допустимых $k,n$? Объясните.
\item Известно, что $-0{,}9<u<-0{,}8$ и $-0{,}9<v<-0{,}8$. Оцените $u/v$.
\item Решите два неравенства и запишите ответы промежутками:
\[
\text{а)}\ -2x>14;\qquad
\text{б)}\ 2x+3\ge5x-6.
\]
\item Решите неравенство $3(x-1)+5>7-3x+2$.
\item Укажите множество решений каждого неравенства:
\[
\text{а)}\ x+6>x;\qquad
\text{б)}\ x+6<x.
\]
\item Найдите все \emph{отрицательные целые} решения неравенства
\[
\frac{13x-1}{15}-\frac{2x-1}{5}<x-\frac{x-2}{3}.
\]
\item Один тюльпан стоит 45 руб., корзинка --- 150 руб. Какое \emph{наибольшее целое количество} тюльпанов можно купить вместе с корзинкой за 2000 руб.?
\end{enumerate}
\vfill
\textcolor{muted}{\small Самопроверка: для деления проверьте знаменатель; при делении неравенства на отрицательное число --- направление знака; для задач на количество --- целочисленность ответа.}
\clearpage

\section*{06 \quad Ответы для самопроверки}
\renewcommand{\arraystretch}{1.65}
\begin{tabularx}{\linewidth}{@{}p{1.05cm}X@{}}
\toprule
\textbf{№} & \textbf{Ответ и ключевой контрольный шаг}\\
\midrule
1 & $0{,}343<a^3<1{,}331$. Куб сохраняет порядок.\\
2 & $4<\dfrac{24}{t}<6$. Больший положительный знаменатель даёт меньшее частное.\\
3 & $-3\le k+n\le4$; \quad $-5\le k-n\le2$.\\
4 & $-6\le kn\le3$. Для $k/n$ конечных общих границ нет: $n=0$ допустимо, а при $k=1$ и $n\to0$ частное неограниченно растёт по модулю.\\
5 & $\dfrac89<\dfrac uv<\dfrac98$.\\
6 & а) $x\in(-\infty;-7)$; \quad б) $x\in(-\infty;3]$.\\
7 & $x>\dfrac76$, то есть $x\in\bigl(\frac76;+\infty\bigr)$.\\
8 & а) $\mathbb R$ (любое действительное $x$); \quad б) $\varnothing$.\\
9 & $x>-\dfrac83$; отрицательные целые решения: $-2$, $-1$.\\
10 & $41$ тюльпан. $45\cdot41+150=1995\le2000$, а $45\cdot42+150=2040>2000$.\\
\bottomrule
\end{tabularx}
\vspace{11mm}
\key{Проверка понимания:} почему перестановка границ при делении зависит от знака знаменателя? Что меняется в ответе, если требуется \emph{целое} число объектов? Умеете ли Вы распознать неравенство, верное при любом $x$, и неравенство без решений?

\vfill
\begin{center}
\small\color{muted} Лёвин Артём Александрович эксклюзивно для Екатерины Скелли
\end{center}
\end{document}
